% ---------------------------------------------------------------------
\section{Additional results}\label{sec:results-supp}
% ---------------------------------------------------------------------
% Tables and figures are read from \analysisdir (default ../../analysis) or from
% local copies in tables/ and figures/. Every number in the text is listed with
% its source cell in paper/results_number_sources.md.

\TODO{The authors check every reported value against the analysis outputs and revise the wording of this section}.

The tables of this section are generated by the analysis pipeline.
Its files T2, T3, T3b, T3c, T3d, T5, T4, T7, T6, TS1, TS2 and TS3 appear as Tables~\ref{tab:t2} to~\ref{tab:ts3} in this order.
Values are means of the seed-averaged policies with the range over the two training seeds in parentheses, and \dag{} marks deterministic methods with a single run.
The example-day figure, the planned ablations and sensitivity studies other than the stress rating were not run.

\subsection{Learning curves and zero-shot transfer}\label{app:res-tables}

\Fig{fig:curves} shows the training curves of the learned methods, with bands over the two seeds.
\Tab{tab:t2} lists the zero-shot results behind \Fig{M-fig:zeroshot} for every station size, including the transfer gap to the reference trained at $N=50$.
That reference is in distribution only at $N=50$, so at $N=100$ and $N=200$ the transfer gap compares two extrapolations.

\begin{figure}[t]
\centering
\resultfig{figS1_training_curves.pdf}{Training curves of the learned methods.}
\caption{Training curves of the learned methods versus environment steps; bands show the range over the two training seeds.}
\label{fig:curves}
\end{figure}

\restable{T2}{\placeholdertable{tab:t2}{Zero-shot transfer across station sizes $N$}{T2.tex}}

\subsection{Constraint satisfaction}\label{app:res-constraints}

Pooled over the test grid, \Tab{tab:t3} gives E4, energy not served, overload, the shortfall relative to the deficit bound of \Eq{eq:deficit} and the capped and admission-reduced requests.
The deficit-bound ratio covers only the methods with the full safety layer, because the variant without departure bound has no such bound.
Per station size, \Tab{tab:t3b} gives E4, the infeasible steps by fallback case and the binding rates of the aggregate limits and of the departure lower bound.
\Tab{tab:t3c} and \Tab{tab:t3d} give event counts summed over all episodes of both seeds, overall and per station size.
Their last column is the episode-level check of Theorem~\ref{M-thm:invariance}(b): unserved EVs in episodes without any infeasible step.
It is zero for every non-conformal method with the departure bound, whereas the variant without departure bound had 1,551 such EVs.
Over all episodes, however, that variant left fewer EVs unserved than \method{}, 1,551 against 6,158 (\Sec{M-sec:res-bound}).
For the conformal variants, the check also counts EVs whose re-certified requirement lay below the shared one or that left before the predicted bound, cases that the theorem does not cover.

\restable{T3}{\placeholdertable{tab:t3}{Constraint satisfaction per method over all test-grid configurations}{T3.tex}}
\restable{T3b}{\placeholdertable{tab:t3b}{E4, infeasible steps by fallback case and binding rates per station size $N$}{T3b.tex}}
\restable{T3c}{\placeholdertable{tab:t3c}{Event counts per method over all evaluated episodes}{T3c.tex}}
\restable{T3d}{\placeholdertable{tab:t3d}{Event counts per method and station size $N$}{T3d.tex}}

\subsection{Cost of the guarantee and constraint handling}\label{app:res-ablation}

\Tab{tab:t5} gives the cost of the guarantee behind \Fig{M-fig:cost}, in $J^{\sm{pen}}$ and in profit, for the attention and MLP pairs.
\Fig{fig:ablation} and \Tab{tab:t4} compare the architectures and constraint-handling variants at the training sizes and in zero-shot transfer.
At $N\ge50$, \method{} had a gap of 0.901, Deep Sets+safety layer 0.674 and Attention+safety layer without departure bound 0.770.
The penalty-trained policy with the layer at deployment reached 0.910, Attention+penalty 0.996 and Attention+PPO-Lagrangian 1.138.

\restable{T5}{\placeholdertable{tab:t5}{Cost of the guarantee per station size $N$}{T5.tex}}

\begin{figure}[t]
\centering
\resultfig{figS2_ablation.pdf}{Gap of the architecture and constraint-handling variants.}
\caption{Gap of the architecture and constraint-handling variants of \Tab{tab:t4} at the training sizes and in zero-shot transfer.}
\label{fig:ablation}
\end{figure}

\restable{T4}{\placeholdertable{tab:t4}{Architecture and constraint handling}{T4.tex}}

\subsection{Conformal departure bound}\label{app:res-conformal}

\Fig{fig:conformal} and \Tab{tab:t7} show coverage, shortfall against the shared requirement, gap and charged energy of \method{}-CP against the miscoverage level $\alpha$.
The achieved coverage was 0.992, 0.950, 0.902 and 0.799 for the targets 0.99, 0.95, 0.90 and 0.80.
Summed over all episodes, the shortfall was 771,853 kWh at $\alpha=0.01$, 12,507 kWh at $\alpha=0.05$ and 2,568 kWh with known departure steps.
For comparison, the uncalibrated point prediction reached a coverage of 0.899 and a gap of 1.322.

\begin{figure}[t]
\centering
\resultfig{figS3_conformal_alpha.pdf}{Coverage, shortfall, gap and charged energy versus $\alpha$.}
\caption{Trade-off of the conformal departure bound: coverage against its target, shortfall against the shared requirement, gap and charged energy per miscoverage level.}
\label{fig:conformal}
\end{figure}

\restable{T7}{\placeholdertable{tab:t7}{Conformal departure bound}{T7.tex}}

\subsection{Computation time}\label{app:res-time}

\Fig{fig:inference} and \Tab{tab:t6} report the inference time per control step and the training time per run.
One training run of \method{} took 1.15 h, and one control step took 1.68 ms at $N=10$ and 7.76 ms at $N=200$.
At $N=200$, Deep Sets+safety layer needed 5.99 ms per step, the rules with the layer at most 2.72 ms and eMPC-HiGHS+safety layer 14.13 ms.

\begin{figure}[t]
\centering
\resultfig{figS4_inference_time.pdf}{Inference time per control step versus $N$.}
\caption{Inference time per control step, policy and safety layer, versus $N$.}
\label{fig:inference}
\end{figure}

\restable{T6}{\placeholdertable{tab:t6}{Computation}{T6.tex}}

\subsection{Relative gaps, stress level and literature configuration}\label{app:res-sensitivity}

\Tab{tab:ts1} reports relative gaps as ratios of sums for the soft-penalty optimum and for both hard optima.
Because $J^{\sm{pen,LP}}$ is often close to zero or negative, the relative $J^{\sm{pen}}$ gap of \method{} exceeded 100\% at $N=5$ and $N=10$, and the main text therefore uses the absolute gap.
\Tab{tab:ts2} covers the stress level $\rho_{\sm{tr}}=3$ kW at $N=20$ and $N=200$.
There, E4 of \method{} was 55.30\% and 65.91\%, against 65.14\% and 73.25\% for Deep Sets+safety layer.
On EV2Gym's 25-charger configuration, Attention+safety layer without departure bound had a gap of 0.363, eMPC-HiGHS with the layer 0.456 and \method{} 0.507 (\Tab{tab:ts3}).

\restable{TS1}{\placeholdertable{tab:ts1}{Relative gaps as ratios of sums}{TS1.tex}}
\restable{TS2}{\placeholdertable{tab:ts2}{Stress level $\rho_{\sm{tr}}=3$ kW}{TS2.tex}}
\restable{TS3}{\placeholdertable{tab:ts3}{Results on EV2Gym's 25-charger literature configuration}{TS3.tex}}

\subsection{Pooled constraint metrics}\label{app:res-example}

\Fig{fig:constraintspooled} shows energy not served and overload energy of every method and transformer rating, pooled over all station sizes of the test grid.
Overload energy appears only for the methods without the safety layer, in line with the zero overload steps of every method with the layer in \Tab{tab:t3c}.
The conformal variants are not shown, and \Tab{tab:t7} gives their shortfall.

\begin{figure}[t]
\centering
\resultfig{figS6_constraints_pooled.pdf}{Constraint metrics per method, pooled over station sizes.}
\caption{Constraint metrics per method and transformer rating, pooled over station sizes.}
\label{fig:constraintspooled}
\end{figure}
